% SPDX-License-Identifier: BSD-3-Clause
\documentclass[10pt,a4paper]{article}
\usepackage[T1]{fontenc}
\usepackage{lmodern}
\usepackage[margin=26mm,headheight=14pt]{geometry}
\usepackage{amsmath,amssymb,amsthm,booktabs,array,tabularx,longtable}
\usepackage{microtype,enumitem,fancyhdr}
\usepackage[hidelinks]{hyperref}
\newtheorem{theorem}{Theorem}
\newtheorem{proposition}[theorem]{Proposition}
\urlstyle{tt}
\newcommand{\code}[1]{\nolinkurl{#1}}
\newcommand{\source}[1]{\code{#1.thy}}
\newcommand{\Nat}{\mathbb{N}}
\newcommand{\cat}{\mathbin{+\!+}}
\newcolumntype{Y}{>{\raggedright\arraybackslash}X}
\setlength{\parindent}{0pt}
\setlength{\parskip}{4pt}
\setlength{\emergencystretch}{2em}
\setlist[itemize]{leftmargin=1.4em,itemsep=2pt,topsep=3pt,parsep=0pt}
\setlist[enumerate]{leftmargin=1.5em,itemsep=2pt,topsep=3pt,parsep=0pt}
\renewcommand{\arraystretch}{1.13}
\pagestyle{fancy}
\fancyhf{}
\fancyhead[L]{\small Evidence Binding Composition}
\fancyhead[R]{\small Technical proof document}
\fancyfoot[C]{\thepage}
\renewcommand{\headrulewidth}{0.3pt}
\raggedbottom
\title{Evidence Binding Composition\\[2mm]\large Finite Allocations, Observable Decisions, and Recovery}
\author{Oraclizer}
\date{8 September 2026}

\begin{document}
\maketitle
\thispagestyle{plain}
\begin{abstract}
This technical note describes the Isabelle/HOL session
\code{Evidence_Binding_Composition}. Finite funding tables are connected to
permitted source-aware operations, fresh completed calls, and recovery of the
same recorded executions. A projection of historical snapshots preserves the
declared call languages, while concrete observers delimit information that
cannot be erased. Financial and regulatory applications consume these
connections. The results concern a conditional mathematical model; physical
source control, current policy supply, durable storage, distributed agreement
and compiled consumers remain implementation obligations.
\end{abstract}

\section{Scope and inherited operations}
The session extends \code{Evidence_Atomic_Binding}, which in turn uses the
existing reservation, message and regulatory models. Throughout this note,
an \emph{actual operation} is an existing HOL transition or reader invoked
by a constructed execution. The word does not assert a physical deployment.
The source statements supply the exact types and assumptions; the propositions
below present their mathematical content in compact notation.

The parent provides immutable source bindings, source-attestation premises,
credited-root funding, current-use checks, confirmed terminal references,
primary publication, and the durable-call and recovery functions. In
particular, it already relates a credited root's holder funding to the replies
of its descendant consumer over common current contexts and requests. It also
already checks candidate recovery data against the complete current source.

The child supplies a finite allocation characterization and connects the
resulting operations to complete call executions. It constructs a historical
projection, proves its primitive update laws, and carries those laws through
the original recovery checks. Two actual applications and several separating
executions show where the connections apply and where they stop.

\paragraph{Notation.}
For a reservation machine $m$, write $J(m)$ for its journal and $P_m(a)$ for
the pooled destination balance of account $a$. A root $r$ has immutable key
$k(r)$ and holder account $a(r,h)$, which includes destination, asset and holder.
Write $F_m(r,h)$ for funding attributed to $k(r)$ at $a(r,h)$. List
concatenation is $u\cat v$; all execution words in this note are finite.

Fixed-margin integer tables and their moves are established mathematical
objects~\cite{diaconisSturmfels1998}. Explicit state mappings and local proof
obligations for behavioral correspondence are also standard~\cite{abadiLamport1991}.
The contribution here concerns their consumption by these specified protocol
operations, not a new general transportation or refinement-mapping theory.

\clearpage
\section{Finite funding margins and constructive realization}
Fix finite lists $R$ of roots and $H$ of holders. Root keys in $R$ are distinct,
holders in $H$ are distinct, every selected root is already credited, and $m$
satisfies the parent reservation contract. A target $T:R\times H\to\Nat$
specifies the selected funding cells. For any complete account $a$, define
\[
 C_T(a)=\sum_{r\in R}\ \sum_{\substack{h\in H\\a(r,h)=a}} T(r,h).
\]
This compares actual accounts. Identical holder numbers on different assets
or destinations do not form one column. Define arithmetic admissibility by
\begin{align}
 \sum_{h\in H}T(r,h)&=\sum_{h\in H}F_m(r,h) &&(r\in R),\label{eq:rows}\\
 C_T(a)&=C_{F_m}(a) &&(\text{every account }a).\label{eq:accounts}
\end{align}
Neither condition calls the execution function or asks whether the target is
safe. Successful execution appears on the other side of the characterization.

\begin{theorem}[Finite allocation characterization]
Under the assumptions above, equations~\eqref{eq:rows}--\eqref{eq:accounts}
hold if and only if a finite supported word of positive canonical ordinary
descendant effects succeeds, reaches $T$ on the selected cells, and preserves
the final pooled balances and the non-lineage reservation state.
\end{theorem}

The exact statement in \source{Funding_Characterization} is:
\par\code{finite_funding_admissibility_iff_realization}.\par
``Supported'' restricts every effect's
root, sender and recipient to $R$ and $H$. Canonical effects use the constructor
in \source{Funding_Realization}; the theorem does not characterize arbitrary
operation grammars or arbitrary fixed permission policies.

\paragraph{Forward proof.}
For each nonempty holder list, choose its first holder as an anchor. Move each
other holder's existing root funding to the anchor, omitting zero quantities.
Then distribute each non-anchor target amount. Equation~\eqref{eq:rows}
ensures that the remaining anchor amount is its target. At every positive move,
the construction supplies the existing current authorization, exact spend
permission and ACTIVE metadata. The amount is bounded by current root funding;
the parent accounting contract supplies the pooled-balance bound as well.

Distinct root keys preserve the rows already processed, even when roots share
destination accounts. Funding outside the selected cells remains unchanged.
The account-total condition, together with the invariant pooled residual,
therefore restores every final pooled balance. Empty root or holder selections
use the empty word.

\paragraph{Reverse proof.}
An accepted effect subtracts and adds the same quantity within one selected
root, so its row sum is unchanged. The supported word preserves outside
funding and the residual $P_m(a)-C_{F_m}(a)$. Equality of final pooled balances
then forces~\eqref{eq:accounts}. The target equality forces~\eqref{eq:rows}.

\paragraph{What is retained.}
The non-lineage comparison is \code{erase_lineage_state}: it clears the funding
map and descendant-history field from reservation state. It does not remove
the real history. For the constructed effect list $w$,
\[
 J(m')=J(m)\cat\operatorname{map}(\mathrm{DescendantEvent},w).
\]
No new source debit, credit or return is introduced. Intermediate pooled
balances may change. No minimal number of transfers or optimal cost is claimed.

\clearpage
\section{An explicit historical projection and its update law}
A historical snapshot stores more than its query functions read. The map
$q$ changes only six fields inside historical snapshots' financial records:

\begin{center}
\begin{tabularx}{\textwidth}{@{}lY@{}}
\toprule
Field & Normalized value\\
\midrule
\code{reservation_at} & No reservation at any key\\
\code{asset_version} & Zero at every asset\\
\code{issued_certificates} & Empty list\\
\code{received_messages} & Empty message state\\
\code{lawful_descendants} & Empty list\\
\code{reservation_clock} & Zero\\
\bottomrule
\end{tabularx}
\end{center}

The current core and endpoint caches are untouched. Historical balances,
application values, owner information, source-effect lists, regulatory state,
terminal records, publication markers and snapshot revisions are retained.
Actual source state and raw journals are retained too. Thus the projection is
not an instruction to delete a production record.

For a callback $e$, let $e(x,s)=(s',r)$, with input $x$ and reply $r$.
\source{Historical_Decision_Projection} proves idempotence and the equations
\begin{align}
 q(q(s))&=q(s),\label{eq:idem}\\
 \operatorname{snd}(e(x,q(s)))&=\operatorname{snd}(e(x,s)),\label{eq:reply}\\
 q(\operatorname{fst}(e(x,q(s))))&=q(\operatorname{fst}(e(x,s))).\label{eq:update}
\end{align}
The reduced callback runs the original function and normalizes its result:
\[
\widehat e(x,s)=(q(\operatorname{fst}(e(x,s))),\operatorname{snd}(e(x,s))).
\]
Its data does not contain a table of future answers or a precomputed safety
decision.

\paragraph{Primitive proof obligations.}
The stored query reader uses balances, application values, regulatory state and
terminal records. Readiness additionally uses source effects, publication and
owner information. None reads the six normalized fields. Case analysis gives
their value and readiness equalities. Current-cache validation and current
authorization read the unchanged current state. Storing a new core result
commutes with normalizing the historical map after that insertion.

These facts cover every observed command and environment constructor. The
sourced callbacks additionally check source availability, source debit/return
gaps and source-coupling guards; these inputs are unchanged. Their update and
reply equations follow from the same primitive checks and actual callbacks.

\begin{theorem}[Finite call projection]
Lift $q$ to the authority state and genesis of the call machine, leaving its
call records intact; denote that lift by $Q$. For every finite client and
environment word $u$,
\[
 \widehat{\operatorname{run}}(u,Q(M))
   =Q(\operatorname{run}(u,M)).
\]
The invocation table, execution results, completions, local response records
and complete call history agree.
\end{theorem}

The proof supplies equations~\eqref{eq:reply}--\eqref{eq:update} to the two
actual durable-call protocols, checks all call-action constructors, and
inducts over $u$. Completed rejection, Busy and Unavailable are included.
This explicit mapping uses standard local-to-behavior reasoning; it does not
appeal to a general existence theorem for a suitable mapping~\cite{abadiLamport1991}.

\clearpage
\section{Source operations, fresh calls and recovery of existing results}
The allocation construction must pass the child dispatcher as well as the
parent financial operation. For each move, \source{Source_Realization_Link}
installs the specified endpoint context through an environment input. It then
consumes the confirmed terminal record and evidence index, primary-publication
marker, and the current regulatory receiver state. The resulting parent state
and reply equal the actual descendant operation's result, including rejection.
Authenticity of the context input remains an implementation obligation.

\source{Source_Call_Realization} translates every coupling input into the
sourced durable-call language. Source attempts, receipt arrivals and context
changes become authority inputs. A client command first sets the explicit
availability input, completes an actual full-cache refresh, then completes
the unchanged command through \code{Execute_Current}. The refresh installs
the current snapshot; it does not manufacture a successful client result.

\paragraph{Freshness and completion.}
The translator reserves two identifiers per coupling input in a fixed
namespace. A lower bound on unused identifiers advances with the word.
Existing invocation, result and completion entries are preserved. For each
fresh command it reuses the parent's begin, dispatch, collect and complete
program. The generated sourced state projects exactly to the original coupling
execution, and each client position completes that execution's actual reply.

This is an uninterrupted finite construction. Arbitrary future words may
contain crashes, disconnections, lost responses and policy inputs; the
projection theorem preserves their behavior without promising successful
delivery. Durable result reuse itself is established systems practice, as
exemplified by RIFL~\cite{leeEtAl2015}. The additional question here is how the
recorded result is tied to the particular source and financial consumer.

\paragraph{The recovery comparison is preserved.}
For a source machine $M$, the original checker compares a candidate's genesis
and ordered entries with those of $M$ and validates its recorded callback
history. Internal replayability alone is insufficient. Successful recovery
returns the actual authority state. Restored dispatch requires a connected
source and full replica equality before using the existing call dispatcher.

\begin{theorem}[Recovery followed by projected continuation]
Suppose the original source satisfies its generated replay contract, its
current candidate is accepted, and the recovery connection is available.
Restore and dispatch using that original
runtime. For every finite suffix, projected execution has the same completed
reply table and call history as the corresponding original dispatch and suffix.
\end{theorem}

\paragraph{Proof sketch.}
The parent recovery theorem identifies the restored state with the current
authority state. Its guarded-dispatch theorem identifies the resulting source
machine with the original dispatcher result. Applying the finite call
projection to that machine proves the suffix equalities. Earlier invocation
identifiers and completed replies are retained: a fresh machine is never
substituted after the recovery boundary.

\source{Integration_Transport} also decodes ordered lists of entry blocks by
concatenation. Encoding entries as singleton blocks decodes to the original
candidate. A complete partition can therefore feed the unchanged current-source
check. This is a theorem about structured HOL entries, not serialization of
arbitrary function-valued state into finite bytes.

\clearpage
\section{Two applications through the actual completion boundary}
\subsection{Same pooled balances, different financial decisions}
The monetary example begins at the actual sourced genesis. Its prefix applies
two source effects, produces their receipts, records terminal decisions,
credits the destination, publishes primary state and reconciles the reservations.
A lawful transfer establishes the initial lineage table. The constructed
exchange then changes the attribution while restoring pooled balances:
\[
 \begin{array}{c@{\qquad}c}
 \text{lineage seed}&\text{after exchange}\\[1mm]
 \begin{pmatrix}0&5\\5&0\end{pmatrix}&
 \begin{pmatrix}5&0\\0&5\end{pmatrix}
 \end{array}
\]
Rows are roots 17 and 23; columns are holders 3 and 4 of the same destination
asset. Both pre-probe states have pooled balances $(5,5)$.

\begin{center}
\begin{tabularx}{\textwidth}{@{}lccY@{}}
\toprule
Pre-probe state & Root 17 at holder 4 & Amount & Fresh completed reply\\
\midrule
Lineage seed & 5 & 1 & \code{Descendant_Executed}\\
After exchange & 0 & 1 & \code{Request_Rejected}\\
\bottomrule
\end{tabularx}
\end{center}

Both probes request root 17, holder $4\to3$, one ordinary unit, and the same
current permission. The inherited funding threshold supplies the two actual
parent replies. Terminal and current-context facts connect those replies to
the source dispatcher; the fresh-call translation completes them. The longer
prefix uses a later unused identifier, not a fabricated cached result.

\source{Funding_Completion_Separation} then restores each original generated
machine. Retransmission uses its real recorded reply, and every finite
projected continuation retains it. This provides an actual completed-call
distinction from equal pooled balances. Pooled equality is a property of the
states \emph{before} the probes; the successful extra transfer changes balances.

\subsection{Regulatory state is an independent input to consumption}
The regulatory application reuses the parent's monetary credit and separate
regulatory-attestation profile. It generates fresh observed calls from the
original regulatory genesis, reaches a published freeze, rejects ordinary use,
and then executes authorized enforcement on that same call machine.

\begin{center}
\begin{tabularx}{\textwidth}{@{}lYl@{}}
\toprule
Call identifier & Actual request after the freeze & Completed result\\
\midrule
$(92,21)$ & Ordinary descendant transfer & Rejected\\
$(92,23)$ & Authorized enforcement transfer & Executed\\
\bottomrule
\end{tabularx}
\end{center}

The ordinary denial does not consume the funding or remove the freeze.
The later enforcement creates a historical snapshot at revision 12 whose
normalization is nontrivial, while preserving query values and readiness.
Both completion entries coexist on the enforced machine. The original
current candidate restores it; redispatch and every projected finite suffix
retain both results. The exact source theorem is
\code{both_regulatory_decisions_survive_recovery_projection_and_future_calls}
in \source{Regulatory_Composition_Example}.

The regulatory certificate remains a supplied attestation input. This
application does not construct its physical producer. Its purpose is to show
that the call and recovery connection also consumes actual regulatory state,
rather than treating available funding as sufficient authorization.

\clearpage
\section{Separating executions and controls}
\subsection{Current snapshots, raw order and original reachability}
A current-cache input can distinguish complete current snapshots. Inject a
candidate cache and invoke the existing zero-time client. Its response is
an \code{Internal_Completed} effect reply exactly when the cache equals the current snapshot;
otherwise it is Busy. Hence equality of every such probe response is
equivalent to equality of the current snapshots. A revision number alone
cannot replace that comparison. The clock example supplies equal revisions
with different full snapshots and demonstrates the changed response after
removing only the current-cache guard.

Raw observers impose different requirements. The journal reader returns the
entire event list, and the available source-effect reader returns its full
effect list. For successful canonical words $u,v$ from the same start,
\[
 J(\operatorname{run}(u,m))=J(\operatorname{run}(v,m))\quad\Longleftrightarrow\quad u=v.
\]
Cancel the common journal prefix and use injectivity of the descendant-event
constructor. Two successful orders of the example exchange reach the same
selected funding and give different raw journal replies. Thus ordered
concatenation is available, while arbitrary permutation is not justified.

Reply preservation also differs from original-state reachability. Every
original generated observed authority stores its current full capture in the
current historical slot. After one clock tick, $q$ changes the historical
clock to zero while keeping the current clock at one. The projected state
violates that original invariant and is not any original generated authority
state, although its callbacks retain their replies. It belongs to the reduced
representation, not to the original full-state recovery interface.

\subsection{A nonempty fixed policy obstructs an admissible table}
Hold a current context fixed whose spend permissions admit positive even
amounts in both directions for both roots. The actual two-unit request
succeeds, so the policy is not an empty denial policy. Every successful move
changes two funding cells by an even amount; every rejected request leaves
funding unchanged. Induction gives a parity invariant for every finite
descendant-only word under that context.

The target below satisfies the same root and account margins but changes
cell parity, making it unreachable under the fixed policy:
\[
 \begin{pmatrix}0&5\\5&0\end{pmatrix}
 \ \longrightarrow\ 
 \begin{pmatrix}1&4\\4&1\end{pmatrix}.
\]
This does not contradict the allocation constructor, which supplies permitted
current contexts for its chosen moves. No unrestricted fixed-policy
reachability theorem is inferred from the arithmetic margins.

\paragraph{The two-leg control.}
Remove only exact spend-tuple membership from the existing descendant
function. Under the same fixed context, first move one unit of root 17 from
holder 4 to holder 3, then one unit of root 23 from holder 3 to holder 4.
Both modified calls actually return Executed and reach the whole target table,
with final pooled balances restored and no new source effects or credits.
The original consumer rejects the first unit call and also the second at that
intermediate state. The proof derives the two actual replies before calculating
the final table. These are controls on formal functions, not compiled-product
mutation tests. See \source{Even_Amount_Policy_Boundary}.

\clearpage
\section{Implementation boundary and reproduction}
The formal constructions preserve the parent's external obligations. They do
not obtain authenticity by changing a client-supplied Boolean, or infer that an
unknown source effect is absent. The following implementation work remains
separate from the statements summarized here.

\begin{description}[leftmargin=0pt,style=nextline,font=\normalfont\bfseries]
\item[Source and terminal evidence.]
Physical adapters must provide immutable effect identity, deduplication,
persistent no-effect fences and authenticated current outcomes. Signature,
roster and stable-source assumptions must be related to their actual producers.
A timeout or lost acknowledgment does not authorize a refund. A later
compensating action has a new identity; it does not replace a stable same-key
finalized fact with a reversed fact.
\item[Current policy and persistence.]
Context installation is an explicit environment input. Its authenticity and
currentness are not established by that constructor. The model also requires
an atomic durable execution/result record and a complete current recovery
source. Their implementation across failures and restart requires evidence.
\item[Observation and settlement.]
Public getters, raw storage, events and timestamp/revision encodings need their
own correspondence. Primary publication does not establish simultaneous
application at every physical domain. Secondary-progress telemetry does not
perform settlement. A token's replay rejection is not automatically the
formal call kernel's retransmission of an earlier reply.
\item[Compiled and distributed behavior.]
The parent API inventory is a source-declaration inventory, not a
compiler-derived ABI certificate. Encoding, compiler output, deployed consumers,
distributed terminal agreement and linearizability of the complete deployed
API are outside this result. Raw observations are not hidden to imply atomicity.
\end{description}

\subsection*{Reproduction commands}
Use Isabelle2025-2 and the dated AFP ADS\_Functor archive
\url{https://isa-afp.org/release/afp-ADS_Functor-2026-02-06.tar.gz}.
Its SHA-256 is:
\par{\small\code{10d6fa8671c461022ae5e71859a07610d616681fed79e2f1c81029a124203c85}}\par
From the repository root, the dependency path must directly contain its
\code{ROOT}. Source checking and native proof checking are separate commands:

{\small
\begin{verbatim}
node Evidence_Binding_Composition/verify-source.mjs
isabelle build -b -j 1 -o threads=1 -o parallel_proofs=0 \
  -d /path/to/ADS_Functor -d . Evidence_Binding_Composition
\end{verbatim}
}

The session sets \code{document = false}. With \code{pdflatex} and
\code{bibtex} on the path, the standalone document builder uses a separate
output directory:
{\small
\begin{verbatim}
node Evidence_Binding_Composition/build-document.mjs \
  output/Evidence_Binding_Composition
\end{verbatim}
}
The public source manifest records the relevant input files; a reproduction
should identify the repository commit, tool versions and command used.
This note does not assign a publication revision, deployment status or
external assurance designation. The source guide names files within the
session directory in the repository.

\clearpage
\appendix
\section{Source guide}
The table gives the relevant theory files, not a substitute for their typed
statements. Parent results are reused under their original assumptions.

{\small
\begin{longtable}{@{}p{0.43\textwidth}p{0.53\textwidth}@{}}
\toprule
Theory & Role in the connection\\
\midrule
\endhead
\source{Funding_Realization} & Constructive gather/distribute word; successful guards, outside-cell and financial frames.\\
\source{Funding_Characterization} & Arithmetic margins if and only if an actual supported successful allocation image exists.\\
\source{Source_Realization_Link} & Current context and terminal/publication premises; actual source-generated two-root prefix.\\
\source{Source_Call_Realization} & Fresh identifiers, refresh and command completion; source-word correspondence and recovery continuation.\\
\source{Historical_Decision_Projection} & Six-field historical normalization, primitive equations and finite call projection.\\
\source{Integration_Transport} & Lossless ordered entry blocks; original restore/dispatch followed by projected continuation.\\
\source{Integration_Boundaries} & Exact current-cache and raw-reader distinctions; original recovery admission.\\
\source{Integration_Examples} & Actual clock and cache examples; nontrivial normalization and cache-guard control.\\
\source{Historical_Reachability_Boundary} & A reply-preserving projected history outside original generated states.\\
\source{Funding_Completion_Separation} & Equal pooled balances with different fresh completed funding decisions, retained after recovery.\\
\source{Regulatory_Composition_Example} & Actual observed monetary/freeze/rejection/enforcement calls and retained completions.\\
\source{Funding_Order_Boundary} & Successful exchange orders with equal funding and different raw journal replies.\\
\source{Funding_Guard_Controls} & Single exact-spend-guard removal on a source-generated financial state.\\
\source{Even_Amount_Policy_Boundary} & Fixed-policy parity obstruction and both actual forbidden unit moves under the same guard removal.\\
\bottomrule
\end{longtable}
}

\begingroup
\small
\bibliographystyle{plain}
\bibliography{root}
\endgroup
\end{document}
